%
% crash_attribution.tex
%
% A Z specification of display-crash attribution: which scenes the Hub blames
% when the display ends a dial.  The state is a sequence of dials (generations
% of the connection), each with an ack ledger of the scenes written to it and the
% scenes whose acks the Hub has read; a display that installs a scene on
% receipt, holds its ack until a buffer swap has completed, and then sends it;
% a dial ended by the display or by the Hub; a death discovered by one of several
% parties and queued; a carried set of suspects from Hub-ended dials; a charge
% made by the one replicator, first at its next push or idle tick; a per-scene
% strike count with quarantine at the threshold; isolation, one scene in flight
% at a time, the ack being the proof.
%
% This is the FIXED design: the leader's rulings T10, T11, T12, T13 and T14.
%   T10  the suspects of a display-ended dial are the scenes written to it and
%        not acked; each death is charged once, whoever discovers it; a
%        Hub-ended dial charges nothing; an empty set charges nothing.
%   T11  the ack means a frame survived the scene.  The display installs a scene
%        on receipt and holds its ack; after the next completed buffer swap it
%        sends every held ack, in receipt order, whether or not the scene was
%        visible.
%   T12  when a Hub-ended dial settles with a non-empty ledger, that ledger joins
%        a carried suspect set.  Reading any ack on any later dial clears the
%        carried set.  A display-ended death charges its own ledger together
%        with the carried set, then clears it.  Each death is queued at
%        discovery and charged first at the next push or idle tick; the queue
%        has no silent bound.
%   T13  a third retire cause, unresponsive, for a dial the keepalive or the
%        isolation probe gives up on while the display may still be alive.  It
%        queues the dial's suspects as a display death does, and also carries
%        them as a Hub end does, because the display may outlive the dial.
%   T14  a display-ended (or unresponsive) capture is taken by the dial's own
%        reader, after a non-blocking drain of the socket and after delivering
%        every complete frame, and before shut().  The read is two steps, Recv
%        (bytes taken off the socket) and Deliver (frame applied to the ledger),
%        and the capture is guarded on both being empty.
%
% Why this model exists: the recurrence rule.  Rounds 4 to 7 of the final review
% on the lux-ghhf branch each found a path on which a display death is blamed on
% the wrong scene, or on none.  Round 7 found that T10's premise, a display that
% acks every scene except the poison one, is false: the real display acks on
% receipt and paints in a later frame (review-r7-silent.md F1; the evaluation
% eval-gvr-r1.md M1).  T11 alone was then shown unsound by a counterexample
% (a poison scene installed on a dial the Hub retires crashes the display on the
% next dial, and an innocent scene is charged); T12 is the repair.  The code it
% governs: src/punt_lux/domain/hub/{batch_sender,recovery,display_connection,
% dial_generation,dial_socket,dial_drain,dial_reader,display_listener,
% display_replies,listener_thread,reply_sink,settlement,suspect_book,death_queue,
% death_charger,dial_death,liveness_probe,replicator}.py and the display's held ack
% in src/punt_lux/display/client_sender.py (its release in
% src/punt_lux/display/frame_turn.py, its hold in
% src/punt_lux/display/render_loop.py).  ADRs: DES-102 (attribution, isolation and
% quarantine), DES-093 (the Hub never controls the display's process), and the
% companion model listener_lifecycle.tex (a dial's generation, its ledger and its
% single end).
%
% Type-check:  fuzz docs/crash_attribution.tex
% Model-check (deadlock-freedom; and the negated goals A1 to A4, see Verification):
%   probcli docs/crash_attribution.tex -model_check -coverage -p MAX_OPERATIONS 200
%   probcli docs/crash_attribution.tex -model_check -nodead -bf \
%     -p MAX_OPERATIONS 200 -goal "<A1 or A2 or A3 or A4>"
%
% Fidelity controls (each differs from this file in the stated places only):
%   docs/crash_attribution_ack_on_receipt_buggy.tex  the ack is sent at receipt
%                                                    (the code before T11)
%   docs/crash_attribution_no_carry_buggy.tex        T11 without T12
%   docs/crash_attribution_inflight_ack_buggy.tex    capture without the drain
% Design variant:
%   docs/crash_attribution_inflight_ack_drain.tex    the drain made part of the
%                                                    capture, not a guard
% T10-era controls (charge only on the replicator's own redial, charge the next
% scene sent, charge a Hub-ended dial, charge a dial twice): see the coverage
% document for their status against this specification.
%
\documentclass[a4paper,10pt,fleqn]{article}
\usepackage[margin=1in]{geometry}
\usepackage{fuzz}
\usepackage[colorlinks=true,linkcolor=blue,citecolor=blue,urlcolor=blue]{hyperref}


\begin{document}

\title{Crash Attribution: a Z Specification}
\author{Formal model of who is blamed when the display dies}
\date{October 2026}
\maketitle

% ============================================================
\section{Introduction}
% ============================================================

A display that crashes while it renders a scene does so a frame after it has
read the scene, so the write that carried the scene has succeeded, and the
display may long since have acknowledged it.  The Hub learns of the death later,
by whichever of several parties looks first: the reader sees end of file, the
keepalive's probe fails, an MCP tool redials, the replicator's next send
breaks.  The crash-loop quarantine rests on one sentence: \emph{a scene that
kills the display again and again is charged each time, and quarantined at the
threshold.}  Four review rounds in succession found a way for that sentence to
fail, each from a different side.  The first three were about \emph{who finds
the death}; the fourth is about \emph{what the ack means}:

\begin{itemize}
  \item \textbf{Silent F1 (round~6).}\ The charge is made only where the
    replicator's own redial fails; a death healed by another party's redial
    leaves no evidence, and the loop runs for ever.
  \item \textbf{Code F1 (round~6).}\ A write woken by the Hub's own retirement of
    the dial is charged to the scenes in flight; the scenes did nothing.
  \item \textbf{Silent F1 (round~7).}\ The display acks a scene \emph{on receipt}
    and paints it in a later frame.  A scene that kills the display while it
    paints is already acked, the dead dial's ledger is empty, and nothing is
    charged; in batching mode the ledger holds later, innocent scenes instead.
  \item \textbf{The counterexample to T11 alone.}\ If the ack is held until a
    swap, a scene installed on a dial and then \emph{held} survives that dial's
    retirement by the Hub.  When the display crashes painting it, on the next
    dial, that dial's ledger holds only the scenes written there.  The culprit
    is on no live dial's ledger.  (Section~\ref{sec:fidelity}.)
\end{itemize}

The defects share a cause: \emph{the decision whom to charge was made from the
circumstances of the discovery, or from an ack that did not mean what the ledger
needed it to mean, and not from what the display actually held when it died.}
T10 made the dial's own record the evidence; T11 gives the record its meaning;
T12 closes the one gap that a per-dial record leaves, the display that outlives
a dial.

\paragraph{The design, by responsibility.}
\begin{description}
  \item[The display installs, holds, swaps, acks.]  A scene is received and
    installed in one step.  Its ack is held.  When a buffer swap completes,
    every held ack is sent, in receipt order, whether or not the scene was
    visible in that frame (a scene in a hidden tab or a minimised frame is
    acked at the swap like any other).  A scene whose install or first frame
    kills the display is therefore never acked.  The swap is not guaranteed: a
    minimised or throttled display stops completing frames, and held acks pile up.
  \item[A dial owns its ledger.]  Every dial records the scenes written to it and
    the scenes whose acks the Hub has \emph{read}; the \emph{ledger} is the
    difference.  An ack that the display has sent and the Hub has not yet read is
    \emph{in flight}.  An ended dial keeps its ledger: it is the evidence of the
    death.
  \item[A dial has one end, with one cause.]  The display ends it (a crash, a
    deliberate stop) or the Hub does (a drop, a reply skew), or the Hub gives up on a display
    that may still be alive (the keepalive's or the isolation probe's
    \emph{unresponsive} retire, T13).  The cause is recorded when the dial ends
    and never overwritten.
  \item[A death is settled once, by anyone, and queued.]  \emph{Discovery} of an
    ended dial is a step by any discoverer.  The first discovery settles the
    dial.  A display-ended dial's suspects are its ledger together with the
    carried set; they are queued, and the carried set is cleared.  An
    unresponsive dial's suspects are queued in the same way, and then the
    suspects are carried, because the display may outlive the dial.  A Hub-ended
    dial's ledger joins the carried set and charges nothing.  Who discovers makes
    no difference to what is queued (T10).
  \item[The drain.]  The Hub's read of an ack is two steps on the reader's thread:
    the bytes leave the socket (\emph{Recv}), and the frame is applied to the
    ledger (\emph{Deliver}).  A display-ended or unresponsive dial is captured
    only when no ack is in flight on the socket and none is held undelivered.  The
    capture is made by the dial's own reader (a non-reader discoverer signals and
    does not capture, because a dial has one reader): it reads what the socket
    already holds, non-blocking, until it would block or reaches end of file;
    delivers every complete frame; and only then takes the ledger, before
    \texttt{shut()} discards the receive queue.  Without this an ack the display
    sent for a scene it survived is still in the socket or in the reader's own
    buffer when the dial is retired, and the scene is charged although it was
    acked.  Reading the socket to end of file is not enough: the bytes may already
    be in the reader's hands (Section~\ref{sec:fidelity}).
  \item[The carried set.]  The display outlives a Hub-ended dial, and with it the
    scenes it installed there and holds.  The Hub cannot tell which of the
    unacked scenes the display holds, so the whole ledger is carried.  (A dial the Hub gave up on is in the same
    case.)  Reading
    \emph{any} ack on any later dial clears the carried set, because that ack
    follows a swap after the install of every scene the display already held:
    those scenes survived it.
  \item[One charger, charging first.]  The replicator alone takes the queue and
    charges it, once per death, before it filters or sends anything in its next
    push, and on its idle tick.  The queue has no bound that drops a death.  A
    scene at the strike threshold is quarantined: dropped from the dirty set and
    never sent again.
  \item[Isolation, the ack as proof.]  Once any scene has a strike the replicator
    sends one scene at a time and waits for that scene's ack to be \emph{read}
    before the next.  The ack is the proof of survival: a pong is not, since it
    can precede the frame.  A death in isolation therefore charges a ledger of at
    most the one scene in flight, together with the carried set.
\end{description}

\paragraph{Bounds the design accepts, and the model's bounds.}
These are properties of what attribution can know, stated so that the code and
the ADR state the same ones.
\begin{description}
  \item[B1  A scene shown later is acked before it is painted.]  A scene the user
    brings into view after its swap (a background tab, a restored minimised
    frame) was acked at the swap.  If its first paint kills the display, the
    scene is in no ledger and the death is not attributed to it.  The model
    cannot show this as a trace, because the poison scene never completes a swap
    (below); it is a bound of T11, not of the model.
  \item[B2  Over-charge: the suspects are a superset.]  The ledger holds every
    scene written to the dial and not acked, including scenes still in transit and
    scenes held but not yet swapped.  A death charges them all.  In batching mode
    an innocent scene written alongside the culprit takes the same strike; in
    isolation a ledger is at most one scene written after the first strike, plus
    whatever batching left unacked.  The carried set adds the ledgers of every
    Hub-ended or unresponsive dial since the last ack read.  An innocent scene is quarantined
    only if it was a suspect at the threshold's worth of deaths.
  \item[B3  The swap stall.]  If the display stops completing swaps, every scene
    received since the last swap is held, unacked, and a later death charges
    all of them.  A2 survives: no scene outside the dial's ledger and the
    carried set is ever charged.  The bound is the number of scenes written to
    that dial since its last ack read, plus the carried set.  In isolation the
    stall blocks the next send, and the isolation probe's timeout retires the dial
    as unresponsive (T13): its ledger is queued for charging, as a display death
    is, and carried, because a stalled display is alive and still holds the
    scenes.  For an unresponsive dial the retire instant is the capture: an ack
    that a live display sends after that point is not read, and its scene is
    charged although the display sent it (in the model, \texttt{Swap} after
    \texttt{Unresponsive} adds nothing to \texttt{flight}).  This is the
    over-charge DES-102 accepts: a scene can be struck when a loaded display
    misses its pings.
  \item[B4  A crash with no dial live is not observed.]  If the display dies
    while the Hub has no live dial, nothing finds the death.  The suspects stay in
    the carried set and are charged at the next display-ended death unless an
    ack read clears them first (the display is then a new one, and the old
    crash is lost).  A1 counts only crashes observed on a live dial.
  \item[O7  Ordering: no earlier dial's scene is installed after a later dial's.]
    This is what makes the carried set's clear sound, and it is a property of
    the code, not a timing assumption.  A scene from a socket that has not
    identified is refused, with no ack (\texttt{render\_loop.py},
    \texttt{hub\_reconciliation.py}).  The new dial's \texttt{ConnectMessage}
    precedes its scenes on the same stream; when the display handles it, it
    closes the predecessor that declared the same \texttt{HubId} (DES-068
    preemption, \texttt{hub\_reconciliation.py}, \texttt{socket\_server.py}),
    discarding every byte the predecessor has not read.  A predecessor already
    removed reached end of file, so all its bytes were read first.  The
    preemption is keyed on \texttt{HubId} on both the same-host and the
    cross-host leg (W11), so the argument holds cross-host.  The residue is
    narrow: a predecessor that has not yet identified, and a \texttt{HubId}
    that changes within one Hub process.  The model installs only on a live
    dial, and keeps the held acks a set: which acks the Hub has read is not
    what the goals observe once the drain has run.
  \item[B7  Several Hubs on one display.]  With two Hubs on one display (DES-089)
    a death caused by Hub~B's scene ends Hub~A's dial with cause display, and Hub~A
    charges its own ledger together with its carried set, although none of its
    scenes caused the death.  Under T11 Hub~A's suspects are everything it sent
    that was received since the last swap, not only scenes in transit; in a
    crash loop driven by Hub~B, Hub~A (isolating, once struck) has its one scene
    in flight struck on every turn, at the rate of Hub~B's culprit.  The bound
    is Hub~A's ledger and carried set per death of the shared display.  The
    model has one Hub; A2 is a statement about a death of a display this Hub was
    attached to, and does not exclude this.
  \item[B6  The poison scene.]  The poison scene kills the display at its
    install or in the first frame after it: it never completes a swap.  Any
    other scene is acked at the swap after its install, and the display may
    also end spontaneously.  Dials: three values (the sentinel, the original and
    its replacement), because a death must be discovered after a redial; the
    runs at three dials and at three scenes in Section~\ref{sec:verify} show the
    result does not depend on the least carrier.
\end{description}

% ============================================================
\section{Constants and Free Types}
% ============================================================

Two dials besides the sentinel: the original and its replacement, the least that
lets a death be discovered after a redial and a carried scene be charged on the
next dial.  Two scenes, one of them the poison scene, the least that lets a scene
be charged for a death it did not cause.  Threshold two: the least for which one
death is not yet a quarantine, so that the loop has a length.  Two discoverers:
the replicator, and any other party.  The traces of the controls
(Section~\ref{sec:fidelity}) show that these bounds reach every defect.

\begin{axdef}
  maxDial : \nat \\
  maxScene : \nat \\
  thr : \nat \\
  poison : \nat
\where
  maxDial = 2 \\
  maxScene = 2 \\
  thr = 2 \\
  poison = 1
\end{axdef}

\begin{zed}
  Sc == 1 \upto maxScene
\end{zed}

\begin{zed}
  Dl == 0 \upto maxDial
\end{zed}

\begin{zed}
  LV ::= lup | ldown
\end{zed}

\begin{zed}
  ENDT ::= enone | edisp | ehub | eunr
\end{zed}

\begin{zed}
  DISC ::= drepl | dother
\end{zed}

\begin{zed}
  PH ::= pidle | ppush
\end{zed}

\begin{axdef}
  bump : 0 \upto thr \fun 0 \upto thr
\where
  bump = \{ 0 \mapsto 1, 1 \mapsto 2, 2 \mapsto 2 \}
\end{axdef}

The function \emph{bump} adds a strike, and stays at the threshold.  Dial~$0$ is
a sentinel: the connection before its first dial, settled from the start.  The
replicator is in one of two phases: idle, or inside a push.

% ============================================================
\section{State}
% ============================================================

The first group is the connection and the display.  \texttt{cur} is the latest
dial, \texttt{up} says whether it is live.  \texttt{ended} records the cause of
each dial's end.  \texttt{written} are the scenes the Hub wrote to each dial,
\texttt{rcvd} those the display has received (installed) from it, \texttt{hold}
the scenes the display holds an ack for, each with the dial it was installed
from, \texttt{flight} the acks sent on the current dial and still in the socket,
\texttt{rbuf} the acks the reader has taken off the socket and not yet applied,
and \texttt{acked} the acks the Hub has applied to the ledger.  A scene is
modelled by its id and written at most once per dial; the code rewrites a scene
on the same dial when it changes, and its ledger is a queue of writes.  The
refinement is sound (a scene id is a suspect while any of its writes is
unacked), but A2 is per \emph{write}: an acked earlier version of a scene does
not clear a later, unacked one.  \texttt{sent} is history: every ack the
display has sent on each dial.

The second group is the Hub's record.  \texttt{carry} is the carried set.
\texttt{dqown} and \texttt{dqcar} are the death queue: for each queued dial, its
ledger and the carried set captured with it.  \texttt{settled} are the dials
whose death a discoverer has processed.  \texttt{dirty} are the scenes the
replicator owes the display.  \texttt{strikes} counts the charges against each
scene.  \texttt{phase} says whether the replicator is inside a push.

\texttt{chargedOn}, \texttt{carriedOn}, \texttt{times}, \texttt{pcrash} and
\texttt{sent} are history, never read by a guard.  \texttt{chargedOn} are the
pairs (scene, dial) charged from the dial's own ledger, \texttt{carriedOn} those
charged from the carried set; \texttt{times} counts how often each dial was
charged; \texttt{pcrash} are the dials on which the display, while the dial was
live, died on the poison scene.  The invariants are statements about them.

\begin{schema}{CA}
  cur : \nat \\
  up : LV \\
  phase : PH \\
  ended : Dl \fun ENDT \\
  written : Dl \fun (\power Sc) \\
  rcvd : Dl \fun (\power Sc) \\
  hold : Sc \pfun Dl \\
  flight : \power Sc \\
  rbuf : \power Sc \\
  sent : Dl \fun (\power Sc) \\
  acked : Dl \fun (\power Sc) \\
  carry : \power Sc \\
  dqown : Dl \pfun (\power Sc) \\
  dqcar : Dl \pfun (\power Sc) \\
  settled : \power Dl \\
  dirty : \power Sc \\
  strikes : Sc \fun (0 \upto thr) \\
  chargedOn : Sc \rel Dl \\
  carriedOn : Sc \rel Dl \\
  times : Dl \fun (0 \upto 2) \\
  pcrash : \power Dl
\where
  cur \leq maxDial
\end{schema}

\section{Initialization}

\begin{schema}{Init}
  CA'
\where
  cur' = 0 \\
  up' = ldown \\
  phase' = pidle \\
  \forall d : Dl \spot ended'~d = enone \\
  \forall d : Dl \spot written'~d = \emptyset \\
  \forall d : Dl \spot rcvd'~d = \emptyset \\
  hold' = \emptyset \\
  flight' = \emptyset \\
  rbuf' = \emptyset \\
  \forall d : Dl \spot sent'~d = \emptyset \\
  \forall d : Dl \spot acked'~d = \emptyset \\
  carry' = \emptyset \\
  dqown' = \emptyset \\
  dqcar' = \emptyset \\
  settled' = \{ 0 \} \\
  dirty' = \emptyset \\
  \forall s : Sc \spot strikes'~s = 0 \\
  chargedOn' = \emptyset \\
  carriedOn' = \emptyset \\
  \forall d : Dl \spot times'~d = 0 \\
  pcrash' = \emptyset
\end{schema}

% ============================================================
\section{Operations}
% ============================================================

\subsection{The replicator's push}

A push is the unit in which the replicator charges: \texttt{PushStart} is
enabled only with the death queue empty, so the queue has been charged first
(\texttt{Charge}, the idle tick, takes one queued death and is enabled whenever
the replicator is idle).  Inside the push the replicator may send, and the dial
may end, be discovered and be redialled by other parties meanwhile; the next
push charges whatever they queued.  \texttt{Send} writes a scene to the live
dial.  The strikes it reads are those of the push's start, since nothing charges
inside a push; in isolation the dial's ledger must be empty first, that is, the
previous scene's ack has been read.  A quarantined scene is never sent.

\begin{schema}{PushStart}
  \Delta CA
\where
  phase = pidle \\
  dqown = \emptyset \\
  cur' = cur \\
  up' = up \\
  phase' = ppush \\
  ended' = ended \\
  written' = written \\
  rcvd' = rcvd \\
  hold' = hold \\
  flight' = flight \\
  rbuf' = rbuf \\
  sent' = sent \\
  acked' = acked \\
  carry' = carry \\
  dqown' = dqown \\
  dqcar' = dqcar \\
  settled' = settled \\
  dirty' = dirty \\
  strikes' = strikes \\
  chargedOn' = chargedOn \\
  carriedOn' = carriedOn \\
  times' = times \\
  pcrash' = pcrash
\end{schema}

\begin{schema}{Send}
  \Delta CA \\
  s? : Sc
\where
  up = lup \\
  phase = ppush \\
  s? \in dirty \\
  strikes~s? < thr \\
  (\exists t : Sc \spot strikes~t > 0) \implies (written~cur \setminus acked~cur) = \emptyset \\
  cur' = cur \\
  up' = up \\
  phase' = phase \\
  ended' = ended \\
  written' = written \oplus \{ cur \mapsto written~cur \cup \{ s? \} \} \\
  rcvd' = rcvd \\
  hold' = hold \\
  flight' = flight \\
  rbuf' = rbuf \\
  sent' = sent \\
  acked' = acked \\
  carry' = carry \\
  dqown' = dqown \\
  dqcar' = dqcar \\
  settled' = settled \\
  dirty' = dirty \setminus \{ s? \} \\
  strikes' = strikes \\
  chargedOn' = chargedOn \\
  carriedOn' = carriedOn \\
  times' = times \\
  pcrash' = pcrash
\end{schema}

\subsection{The display}

\texttt{Install} receives a scene and installs it; the ack is held, and a scene
reinstalled from a later dial replaces its earlier hold.  \texttt{Swap} is a
completed buffer swap: it sends every ack the display holds for the live dial
(into \texttt{flight}) and releases every hold.  It never completes while the
poison scene is installed (B6).  \texttt{Recv} is the reader taking an ack off the socket;
\texttt{Deliver} is the reader applying it: the ledger loses the scene, and the
carried set is cleared.  The reader may do both while the dial is live, and
after the display has ended it or the Hub has given up on it, until a
discoverer settles the dial.

\begin{schema}{Install}
  \Delta CA \\
  s? : Sc
\where
  up = lup \\
  s? \in written~cur \setminus rcvd~cur \\
  cur' = cur \\
  up' = up \\
  phase' = phase \\
  ended' = ended \\
  written' = written \\
  rcvd' = rcvd \oplus \{ cur \mapsto rcvd~cur \cup \{ s? \} \} \\
  hold' = hold \oplus \{ s? \mapsto cur \} \\
  flight' = flight \\
  rbuf' = rbuf \\
  sent' = sent \\
  acked' = acked \\
  carry' = carry \\
  dqown' = dqown \\
  dqcar' = dqcar \\
  settled' = settled \\
  dirty' = dirty \\
  strikes' = strikes \\
  chargedOn' = chargedOn \\
  carriedOn' = carriedOn \\
  times' = times \\
  pcrash' = pcrash
\end{schema}

\begin{schema}{Swap}
  \Delta CA
\where
  hold \neq \emptyset \\
  poison \notin \dom hold \\
  cur' = cur \\
  up' = up \\
  phase' = phase \\
  ended' = ended \\
  written' = written \\
  rcvd' = rcvd \\
  hold' = \emptyset \\
  flight' = flight \cup \{ s : \dom hold | hold~s = cur \land up = lup \spot s \} \\
  rbuf' = rbuf \\
  sent' = sent \oplus \{ cur \mapsto sent~cur \cup \{ s : \dom hold | hold~s = cur \land up = lup \spot s \} \} \\
  acked' = acked \\
  carry' = carry \\
  dqown' = dqown \\
  dqcar' = dqcar \\
  settled' = settled \\
  dirty' = dirty \\
  strikes' = strikes \\
  chargedOn' = chargedOn \\
  carriedOn' = carriedOn \\
  times' = times \\
  pcrash' = pcrash
\end{schema}

\begin{schema}{Recv}
  \Delta CA \\
  s? : Sc
\where
  up = lup \lor (ended~cur \in \{ edisp, eunr \} \land cur \notin settled) \\
  s? \in flight \\
  cur' = cur \\
  up' = up \\
  phase' = phase \\
  ended' = ended \\
  written' = written \\
  rcvd' = rcvd \\
  hold' = hold \\
  flight' = flight \setminus \{ s? \} \\
  rbuf' = rbuf \cup \{ s? \} \\
  sent' = sent \\
  acked' = acked \\
  carry' = carry \\
  dqown' = dqown \\
  dqcar' = dqcar \\
  settled' = settled \\
  dirty' = dirty \\
  strikes' = strikes \\
  chargedOn' = chargedOn \\
  carriedOn' = carriedOn \\
  times' = times \\
  pcrash' = pcrash
\end{schema}

\begin{schema}{Deliver}
  \Delta CA \\
  s? : Sc
\where
  up = lup \lor (ended~cur \in \{ edisp, eunr \} \land cur \notin settled) \\
  s? \in rbuf \\
  cur' = cur \\
  up' = up \\
  phase' = phase \\
  ended' = ended \\
  written' = written \\
  rcvd' = rcvd \\
  hold' = hold \\
  flight' = flight \\
  rbuf' = rbuf \setminus \{ s? \} \\
  sent' = sent \\
  acked' = acked \oplus \{ cur \mapsto acked~cur \cup \{ s? \} \} \\
  carry' = \emptyset \\
  dqown' = dqown \\
  dqcar' = dqcar \\
  settled' = settled \\
  dirty' = dirty \\
  strikes' = strikes \\
  chargedOn' = chargedOn \\
  carriedOn' = carriedOn \\
  times' = times \\
  pcrash' = pcrash
\end{schema}

\subsection{The ends of a dial}

\texttt{DisplayEnd} is a spontaneous end, a crash or a deliberate stop, whatever
the ledger holds, empty included; it takes the display's holds with it.
\texttt{HubEnd} is a retire by the Hub (a drop, a reply skew); the
display lives on with what it holds.  \texttt{Unresponsive} is the keepalive's or
the isolation probe's retire of a dial whose display may still be alive (T13): the
display keeps what it holds.  \texttt{CrashInstall} and
\texttt{CrashFrame} are the poison scene's two ways of killing the display: at
its install, or in the first frame after it, which may be on a later dial than
the one that carried it (the counterexample to T11 alone).
\texttt{DisplayEndDown} is a crash while no dial is live (B4).  None of these
charges anything: a death is charged when it is discovered, from the dial that
died.

\begin{schema}{DisplayEnd}
  \Delta CA
\where
  up = lup \\
  cur' = cur \\
  up' = ldown \\
  phase' = phase \\
  ended' = ended \oplus \{ cur \mapsto edisp \} \\
  written' = written \\
  rcvd' = rcvd \\
  hold' = \emptyset \\
  flight' = flight \\
  rbuf' = rbuf \\
  sent' = sent \\
  acked' = acked \\
  carry' = carry \\
  dqown' = dqown \\
  dqcar' = dqcar \\
  settled' = settled \\
  dirty' = dirty \\
  strikes' = strikes \\
  chargedOn' = chargedOn \\
  carriedOn' = carriedOn \\
  times' = times \\
  pcrash' = pcrash
\end{schema}

\begin{schema}{HubEnd}
  \Delta CA
\where
  up = lup \\
  cur' = cur \\
  up' = ldown \\
  phase' = phase \\
  ended' = ended \oplus \{ cur \mapsto ehub \} \\
  written' = written \\
  rcvd' = rcvd \\
  hold' = hold \\
  flight' = flight \\
  rbuf' = rbuf \\
  sent' = sent \\
  acked' = acked \\
  carry' = carry \\
  dqown' = dqown \\
  dqcar' = dqcar \\
  settled' = settled \\
  dirty' = dirty \\
  strikes' = strikes \\
  chargedOn' = chargedOn \\
  carriedOn' = carriedOn \\
  times' = times \\
  pcrash' = pcrash
\end{schema}

\begin{schema}{Unresponsive}
  \Delta CA
\where
  up = lup \\
  cur' = cur \\
  up' = ldown \\
  phase' = phase \\
  ended' = ended \oplus \{ cur \mapsto eunr \} \\
  written' = written \\
  rcvd' = rcvd \\
  hold' = hold \\
  flight' = flight \\
  rbuf' = rbuf \\
  sent' = sent \\
  acked' = acked \\
  carry' = carry \\
  dqown' = dqown \\
  dqcar' = dqcar \\
  settled' = settled \\
  dirty' = dirty \\
  strikes' = strikes \\
  chargedOn' = chargedOn \\
  carriedOn' = carriedOn \\
  times' = times \\
  pcrash' = pcrash
\end{schema}

\begin{schema}{CrashInstall}
  \Delta CA
\where
  up = lup \\
  poison \in written~cur \setminus rcvd~cur \\
  cur' = cur \\
  up' = ldown \\
  phase' = phase \\
  ended' = ended \oplus \{ cur \mapsto edisp \} \\
  written' = written \\
  rcvd' = rcvd \\
  hold' = \emptyset \\
  flight' = flight \\
  rbuf' = rbuf \\
  sent' = sent \\
  acked' = acked \\
  carry' = carry \\
  dqown' = dqown \\
  dqcar' = dqcar \\
  settled' = settled \\
  dirty' = dirty \\
  strikes' = strikes \\
  chargedOn' = chargedOn \\
  carriedOn' = carriedOn \\
  times' = times \\
  pcrash' = pcrash \cup \{ cur \}
\end{schema}

\begin{schema}{CrashFrame}
  \Delta CA
\where
  up = lup \\
  poison \in \dom hold \\
  cur' = cur \\
  up' = ldown \\
  phase' = phase \\
  ended' = ended \oplus \{ cur \mapsto edisp \} \\
  written' = written \\
  rcvd' = rcvd \\
  hold' = \emptyset \\
  flight' = flight \\
  rbuf' = rbuf \\
  sent' = sent \\
  acked' = acked \\
  carry' = carry \\
  dqown' = dqown \\
  dqcar' = dqcar \\
  settled' = settled \\
  dirty' = dirty \\
  strikes' = strikes \\
  chargedOn' = chargedOn \\
  carriedOn' = carriedOn \\
  times' = times \\
  pcrash' = pcrash \cup \{ cur \}
\end{schema}

\begin{schema}{DisplayEndDown}
  \Delta CA
\where
  up = ldown \\
  hold \neq \emptyset \\
  cur' = cur \\
  up' = up \\
  phase' = phase \\
  ended' = ended \\
  written' = written \\
  rcvd' = rcvd \\
  hold' = \emptyset \\
  flight' = flight \\
  rbuf' = rbuf \\
  sent' = sent \\
  acked' = acked \\
  carry' = carry \\
  dqown' = dqown \\
  dqcar' = dqcar \\
  settled' = settled \\
  dirty' = dirty \\
  strikes' = strikes \\
  chargedOn' = chargedOn \\
  carriedOn' = carriedOn \\
  times' = times \\
  pcrash' = pcrash
\end{schema}

\subsection{Discovery and the queue}

A discoverer, whoever it is, finds the current dial ended and unsettled
(\texttt{d?} is an input that no guard reads: that is T10).  If the display ended
it, the drain has run (no ack is in the socket or in the reader's buffer: the
dial's reader has taken and delivered everything), and the ledger together with
the carried set is non-empty, \texttt{DiscoverCharge} queues them and clears the
carried set; if both are empty, \texttt{DiscoverQuiet} settles the dial and
nothing is queued.  \texttt{DiscoverUnrCharge} and \texttt{DiscoverUnrQuiet} do the
same for an unresponsive dial, with one difference: the suspects that were queued
become the carried set.  If the Hub ended the dial, \texttt{DiscoverHub} adds the
ledger to the carried set and queues nothing.

\begin{schema}{DiscoverCharge}
  \Delta CA \\
  d? : DISC
\where
  up = ldown \\
  cur \notin settled \\
  ended~cur = edisp \\
  flight = \emptyset \\
  rbuf = \emptyset \\
  (written~cur \setminus acked~cur) \cup carry \neq \emptyset \\
  cur' = cur \\
  up' = up \\
  phase' = phase \\
  ended' = ended \\
  written' = written \\
  rcvd' = rcvd \\
  hold' = hold \\
  flight' = flight \\
  rbuf' = rbuf \\
  sent' = sent \\
  acked' = acked \\
  carry' = \emptyset \\
  dqown' = dqown \cup \{ cur \mapsto (written~cur \setminus acked~cur) \} \\
  dqcar' = dqcar \cup \{ cur \mapsto carry \} \\
  settled' = settled \cup \{ cur \} \\
  dirty' = dirty \\
  strikes' = strikes \\
  chargedOn' = chargedOn \\
  carriedOn' = carriedOn \\
  times' = times \\
  pcrash' = pcrash
\end{schema}

\begin{schema}{DiscoverQuiet}
  \Delta CA \\
  d? : DISC
\where
  up = ldown \\
  cur \notin settled \\
  ended~cur = edisp \\
  flight = \emptyset \\
  rbuf = \emptyset \\
  (written~cur \setminus acked~cur) \cup carry = \emptyset \\
  cur' = cur \\
  up' = up \\
  phase' = phase \\
  ended' = ended \\
  written' = written \\
  rcvd' = rcvd \\
  hold' = hold \\
  flight' = flight \\
  rbuf' = rbuf \\
  sent' = sent \\
  acked' = acked \\
  carry' = carry \\
  dqown' = dqown \\
  dqcar' = dqcar \\
  settled' = settled \cup \{ cur \} \\
  dirty' = dirty \\
  strikes' = strikes \\
  chargedOn' = chargedOn \\
  carriedOn' = carriedOn \\
  times' = times \\
  pcrash' = pcrash
\end{schema}

\begin{schema}{DiscoverUnrCharge}
  \Delta CA \\
  d? : DISC
\where
  up = ldown \\
  cur \notin settled \\
  ended~cur = eunr \\
  flight = \emptyset \\
  rbuf = \emptyset \\
  (written~cur \setminus acked~cur) \cup carry \neq \emptyset \\
  cur' = cur \\
  up' = up \\
  phase' = phase \\
  ended' = ended \\
  written' = written \\
  rcvd' = rcvd \\
  hold' = hold \\
  flight' = flight \\
  rbuf' = rbuf \\
  sent' = sent \\
  acked' = acked \\
  carry' = (written~cur \setminus acked~cur) \cup carry \\
  dqown' = dqown \cup \{ cur \mapsto (written~cur \setminus acked~cur) \} \\
  dqcar' = dqcar \cup \{ cur \mapsto carry \} \\
  settled' = settled \cup \{ cur \} \\
  dirty' = dirty \\
  strikes' = strikes \\
  chargedOn' = chargedOn \\
  carriedOn' = carriedOn \\
  times' = times \\
  pcrash' = pcrash
\end{schema}

\begin{schema}{DiscoverUnrQuiet}
  \Delta CA \\
  d? : DISC
\where
  up = ldown \\
  cur \notin settled \\
  ended~cur = eunr \\
  flight = \emptyset \\
  rbuf = \emptyset \\
  (written~cur \setminus acked~cur) \cup carry = \emptyset \\
  cur' = cur \\
  up' = up \\
  phase' = phase \\
  ended' = ended \\
  written' = written \\
  rcvd' = rcvd \\
  hold' = hold \\
  flight' = flight \\
  rbuf' = rbuf \\
  sent' = sent \\
  acked' = acked \\
  carry' = carry \\
  dqown' = dqown \\
  dqcar' = dqcar \\
  settled' = settled \cup \{ cur \} \\
  dirty' = dirty \\
  strikes' = strikes \\
  chargedOn' = chargedOn \\
  carriedOn' = carriedOn \\
  times' = times \\
  pcrash' = pcrash
\end{schema}

\begin{schema}{DiscoverHub}
  \Delta CA \\
  d? : DISC
\where
  up = ldown \\
  cur \notin settled \\
  ended~cur = ehub \\
  cur' = cur \\
  up' = up \\
  phase' = phase \\
  ended' = ended \\
  written' = written \\
  rcvd' = rcvd \\
  hold' = hold \\
  flight' = flight \\
  rbuf' = rbuf \\
  sent' = sent \\
  acked' = acked \\
  carry' = carry \cup (written~cur \setminus acked~cur) \\
  dqown' = dqown \\
  dqcar' = dqcar \\
  settled' = settled \cup \{ cur \} \\
  dirty' = dirty \\
  strikes' = strikes \\
  chargedOn' = chargedOn \\
  carriedOn' = carriedOn \\
  times' = times \\
  pcrash' = pcrash
\end{schema}

\subsection{Redial and the charge}

A redial replaces a settled dial.  A redial of an unsettled dial is a discovery
followed by this step; the model keeps them apart because they are two steps in
the code, and any other party's step may come between.  The new dial's scenes
are all marked dirty except the quarantined ones, and the acks in flight on the
old dial are gone with it.  \texttt{Charge} is the replicator taking one queued
death: every scene in its ledger and in its captured carried set takes one
strike, once for the dial, and a scene that reaches the threshold leaves the
dirty set.  \texttt{PushEnd} ends a push.

\begin{schema}{Redial}
  \Delta CA
\where
  up = ldown \\
  cur < maxDial \\
  cur \in settled \\
  cur' = cur + 1 \\
  up' = lup \\
  phase' = phase \\
  ended' = ended \\
  written' = written \\
  rcvd' = rcvd \\
  hold' = hold \\
  flight' = \emptyset \\
  rbuf' = \emptyset \\
  sent' = sent \\
  acked' = acked \\
  carry' = carry \\
  dqown' = dqown \\
  dqcar' = dqcar \\
  settled' = settled \\
  dirty' = \{ s : Sc | strikes~s < thr \} \\
  strikes' = strikes \\
  chargedOn' = chargedOn \\
  carriedOn' = carriedOn \\
  times' = times \\
  pcrash' = pcrash
\end{schema}

\begin{schema}{Charge}
  \Delta CA \\
  g? : Dl
\where
  phase = pidle \\
  g? \in \dom dqown \\
  cur' = cur \\
  up' = up \\
  phase' = phase \\
  ended' = ended \\
  written' = written \\
  rcvd' = rcvd \\
  hold' = hold \\
  flight' = flight \\
  rbuf' = rbuf \\
  sent' = sent \\
  acked' = acked \\
  carry' = carry \\
  dqown' = \{ g? \} \ndres dqown \\
  dqcar' = \{ g? \} \ndres dqcar \\
  settled' = settled \\
  dirty' = dirty \cap \{ s : Sc | strikes'~s < thr \} \\
  strikes' = strikes \oplus \{ s : dqown~g? \cup dqcar~g? \spot s \mapsto bump~(strikes~s) \} \\
  chargedOn' = chargedOn \cup \{ s : dqown~g? \spot s \mapsto g? \} \\
  carriedOn' = carriedOn \cup \{ s : dqcar~g? \spot s \mapsto g? \} \\
  times' = times \oplus \{ g? \mapsto times~g? + 1 \} \\
  pcrash' = pcrash
\end{schema}

\begin{schema}{PushEnd}
  \Delta CA
\where
  phase = ppush \\
  cur' = cur \\
  up' = up \\
  phase' = pidle \\
  ended' = ended \\
  written' = written \\
  rcvd' = rcvd \\
  hold' = hold \\
  flight' = flight \\
  rbuf' = rbuf \\
  sent' = sent \\
  acked' = acked \\
  carry' = carry \\
  dqown' = dqown \\
  dqcar' = dqcar \\
  settled' = settled \\
  dirty' = dirty \\
  strikes' = strikes \\
  chargedOn' = chargedOn \\
  carriedOn' = carriedOn \\
  times' = times \\
  pcrash' = pcrash
\end{schema}

\subsection{Halt}

\begin{schema}{Halt}
  \Xi CA
\where
  cur = maxDial \\
  up = ldown \\
  cur \in settled \\
  phase = pidle \\
  dqown = \emptyset
\end{schema}

% ============================================================
\section{Invariants}
\label{sec:inv}
% ============================================================

Each is a negated goal: the invariant holds if its negation is unreachable
(Section~\ref{sec:verify}).  They are deliberately not conjuncts of the state
schema: in ProB's reading of Z a primed state schema is part of every
operation's guard, so an invariant written there silently disables an operation
whose result would violate it, instead of reporting the violation.  They are
stated over the history variables and the ground truth the model keeps
(\texttt{pcrash}, \texttt{sent}), not over the Hub's own bookkeeping, so that a
design cannot satisfy them by its own definitions.

\paragraph{A1 --- every observed poison death is charged to the poison scene,
whoever discovers it.}
Let a \emph{poison death} be a dial on which the display, while the dial was
live, died on the poison scene (\texttt{pcrash}).  The negated goal:
\[ \exists g \in pcrash \spot g \in settled \land g \notin \dom dqown \land
     (poison \mapsto g) \notin chargedOn \cup carriedOn. \]
A poison death that is settled and no longer queued, and has not charged the
poison scene, is a death lost.  This is the property of the loop, and it is
stronger than reaching quarantine: each poison death takes its strike, so
quarantine follows after $thr$ of them.  It covers a scene that crashes the
install or the first frame after it; discovery by any party, across Hub-ended
dials and while the replicator is away, because the queue keeps a death until
the replicator charges it.  A sharper form, \textbf{A1w}, adds that some scene
\emph{was} charged for the death: the wrong scene took the strike.

\paragraph{A2 --- no scene is charged for a death of a display this Hub was
attached to unless it was in that death.}
The name is narrower than ``a death it was not in'': with several Hubs on one
display a scene of this Hub can be charged for a death another Hub's scene
caused (B7), and A2 does not exclude that.  It is a statement per write (see the
State section).  Two clauses.  Own ledger:
\[ \exists (s,g) \in chargedOn \spot s \notin written~g \lor s \in sent~g. \]
A scene is charged from a dial's ledger only if it was written to that dial and
the display did not send its ack: the ledger is a property of the dial, so no
circumstance of the discovery can put another scene into it, and a scene the
display survived (an ack in flight included) is not a suspect.  Carried set:
\[ \exists (s,g) \in carriedOn \spot \lnot \exists d < g \spot
     ended~d \in \{ ehub, eunr \}
     \land s \in written~d \setminus acked~d \land
     (\forall e : d < e \leq g \spot acked~e = \emptyset) \land
     (\forall e : d < e < g \spot ended~e \neq edisp). \]
A carried scene is charged only if some earlier Hub-ended or unresponsive dial
wrote it without an applied ack (an unresponsive dial re-carries what it charged,
so the earliest such dial is the witness), and no ack has been read since, and no display-ended death in between
has already taken the carried set.  This is the bound of T12 in the model's
terms; an innocent scene inside it is B2's over-charge, not a violation.

\paragraph{A3 --- a Hub-ended dial charges no scene.}  (An unresponsive dial does:
it is the display's end as far as DES-102 is concerned, and A3 says nothing of it.)
\[ \exists (s,g) \in chargedOn \cup carriedOn \spot ended~g = ehub. \]
The guard of \texttt{DiscoverCharge} requires the display to have ended the dial;
a Hub-ended dial's ledger is carried, which is not a charge.  Like A4 this is a
tautology of the fixed design; its content is the guard.

\paragraph{A4 --- each dial's death is charged at most once.}
\[ \exists g \in Dl \spot times~g > 1. \]
\texttt{DiscoverCharge} requires the dial unsettled and settles it, and
\texttt{Charge} removes the death from the queue.  The tautology is by the guard;
two discoverers race for the same dial, and the race is covered
(Section~\ref{sec:verify}).

\paragraph{Deadlock-freedom.}
Every state other than the quiescent one (\texttt{Halt}) has an enabled step: a
live dial can always end, an ended one can always be discovered (the drain is \texttt{Recv} and \texttt{Deliver}, enabled until the dial is
settled), a settled one redialled
unless it is the last, a queued death charged by an idle replicator, a push
ended.

\paragraph{Charge liveness.}
A1 as written is a safety property: a death that stayed queued for ever would
satisfy it.  The model has the liveness property too, by an acyclicity argument.
\texttt{PushStart} needs \texttt{dqown = $\emptyset$}; every other step either
consumes a finite resource (a dial, a scene's turn on a dial, a held ack, a
queued death) or is a push boundary; so no cycle passes through a state with a
queued death and \texttt{Charge} disabled.  With deadlock-freedom checked, every
maximal run reaches \texttt{Charge} for each queued death.  A stalled charger
would break exactly this, which is why the code must keep the replicator's idle
tick charging (O1).

% ============================================================
\section{Fidelity Controls}
\label{sec:fidelity}
% ============================================================

Each control differs from this specification in one place.  Traces are
breadth-first shortest, with the number of states searched.

\paragraph{The ack is sent at receipt ---
\texttt{crash\_attribution\_ack\_on\_receipt\_buggy.tex}.}
The code before T11.  \texttt{Install} sends the ack at once, \texttt{Swap} sends none.
A1 is violated, in 8 steps and 431 states, by \texttt{PushStart; Redial;
Send(1); Install(1); Recv(1); Deliver(1); CrashFrame; DiscoverQuiet}: the scene is
acked, then the display dies painting it, the ledger is empty, nothing is
charged.  This is gvr's trace.  A1w is violated in batching mode, 11 steps and
2657 states, by \texttt{PushStart; Redial; Send(1); Send(2); Install(1); PushEnd;
Recv(1); CrashFrame; Deliver(1); DiscoverCharge; Charge}: the ledger holds the
later, innocent scene 2, and the culprit takes no strike.  A2, A3 and A4 hold:
the ledger is sound; it is the ack that does not mean what the ledger needs.

\paragraph{T11 without T12 ---
\texttt{crash\_attribution\_no\_carry\_buggy.tex}.}
Nothing is carried: \texttt{DiscoverHub} adds nothing to the carried set, and an
unresponsive dial's discovery leaves it alone, so a display-ended death charges
its own ledger only.  Two doors reach the same loss.  Through the Hub's retire,
A1 is violated in 9 steps and 918 states by \texttt{PushStart; Redial; Send(1);
Install(1); HubEnd; DiscoverHub; Redial; CrashFrame; DiscoverQuiet}, and A1w in 12
steps and 6430 states by \texttt{PushStart; Redial; Send(1); Install(1); HubEnd;
DiscoverHub; Redial; Send(2); CrashFrame; PushEnd; DiscoverCharge; Charge}: the
scene is installed and held, the Hub retires the dial, the display crashes
painting it on the next dial, and the innocent scene 2 is charged in its place.
Through the probe's retire (T13), A1w again, 12 steps and 6546 states, with
\texttt{HubEnd; DiscoverHub} replaced by \texttt{Unresponsive;
DiscoverUnrCharge}: the unresponsive dial queues its suspects and carries none,
and the display, still alive, holds the culprit across it.  Removing only the
unresponsive door's carry, \texttt{crash\_attribution\_unresponsive\_no\_carry\_buggy.tex},
reproduces the second trace (A1w) and leaves the first door closed.  A2, A3 and A4 hold.

\paragraph{Capture without the split-read guard ---
\texttt{crash\_attribution\_inflight\_ack\_buggy.tex}.}
\texttt{DiscoverCharge}, \texttt{DiscoverQuiet} and the unresponsive discoveries
keep the guard that no ack is in the socket (\texttt{flight = $\emptyset$}, the
obligation as first worded: ``read the socket to end of file'') and lose the guard
that the reader holds none undelivered (\texttt{rbuf = $\emptyset$}), and capture
\texttt{written $\setminus$ acked}.  A2 (own ledger) is violated, in 10 steps and
2347 states, by \texttt{PushStart; Redial; Send(2); Install(2); Swap; DisplayEnd;
Recv(2); DiscoverCharge(dother); PushEnd; Charge}: the display acked scene 2 and
died; the reader took the ack off the socket and had not yet applied it; a
non-reader retired the dial; the socket is empty, so reading it to end of file
finds nothing, and scene 2 is charged although the display sent its ack and the
Hub received it.  With no guard at all, A2 is violated sooner (9 steps, 1161
states).  The hazard is real, and the guard on the reader's buffer, not the read
of the socket, is what removes it.  A1, A3 and A4 hold.

\paragraph{The drain as part of the capture ---
\texttt{crash\_attribution\_inflight\_ack\_drain.tex}.}
A design variant, not a control.  Discovery carries no guard on the flight or the
buffer; instead each capture takes every ack in the socket and in the reader's
buffer into \texttt{acked} in the same step, and a drain that applies an ack
clears the carried set as any delivery does.  This is what the dial's own reader
does when it drains the socket without blocking, delivers every complete frame,
and only then takes the ledger, before \texttt{shut()}.  A1 to A4 hold, and so
does deadlock-freedom.  The two formulations agree (the state space has the same
size): the model supports the drain, taken by the reader.

\paragraph{Charging a dial the Hub ended, a dial twice, only on the
replicator's own redial, the next scene sent.}
The four controls of the T10 specification, \texttt{crash\_attribution\_hub\_end\_charged\_buggy.tex},
\texttt{crash\_attribution\_double\_charge\_buggy.tex},
\texttt{crash\_attribution\_redial\_only\_buggy.tex} and
\texttt{crash\_attribution\_idle\_stop\_buggy.tex}, are regenerated against this
specification; they reproduce the same defects (A3, A4, A1 and A2), recorded in
the coverage document.  The
\texttt{idle\_stop} control is a substitution, not a one-line mutation: it
replaces the queue-and-charge discovery with the T10-era mechanism that records
a blame and charges the next scene sent.

% ============================================================
\section{Verification}
\label{sec:verify}
% ============================================================

Type-checking is a \texttt{fuzz} pass:
\begin{verbatim}
  fuzz docs/crash_attribution.tex
\end{verbatim}

Deadlock-freedom, with every operation covered:
\begin{verbatim}
  probcli docs/crash_attribution.tex -model_check -coverage -p MAX_OPERATIONS 200
\end{verbatim}
The negated goals, one at a time for a control and as a disjunction for the
design; a pass is ``no counter example, all states visited'':
\begin{verbatim}
  probcli docs/crash_attribution.tex -model_check -nodead -bf \
    -p MAX_OPERATIONS 200 -goal "<A1 or A2 or A3 or A4>"
\end{verbatim}
where the goals, in B syntax, are
\begin{verbatim}
  A1  #g.(g:pcrash & g:settled & g/:dom(dqown)
          & (1|->g)/:chargedOn & (1|->g)/:carriedOn)
  A1w A1 & #s.((s|->g):chargedOn\/carriedOn)           (for the same g)
  A2  #(s,g).((s|->g):chargedOn & (s/:written(g) or s:sent(g)))
      or
      #(s,g).((s|->g):carriedOn & not(#d.(d:0..2 & d<g
          & (ended(d)=ehub or ended(d)=eunr) & s:written(d)-acked(d)
          & !e.(e:0..2 & d<e & e<=g => acked(e)={})
          & !e.(e:0..2 & d<e & e<g => ended(e)/=edisp))))
  A3  #(s,g).((s|->g):(chargedOn\/carriedOn) & ended(g)=ehub)
  A4  #g.(g:0..2 & times(g)>1)
\end{verbatim}
The scenarios are reachable, so the invariants are not vacuous (goal
\emph{found} on the fixed design): a carried scene charged
(\texttt{carriedOn /= \{\}}), the death queue holding two deaths while the
replicator is away (\texttt{card(dom(dqown)) = 2}), the poison scene quarantined
(\texttt{strikes(1) = 2}), an innocent scene quarantined (\texttt{strikes(2) =
2}), an innocent scene charged for a poison death, the poison scene charged through the carried set of an unresponsive dial, and an ack taken off the socket and not yet applied when the display dies (\texttt{rbuf /= \{\}}).  The carried-set clause of A2
is shown to bite by a control that does not clear the carried set on an ack
read, and A1 by a control that bounds the queue at one death and evicts the
oldest.  Results, the bounds and the traces are recorded in
\texttt{crash\_attribution\_coverage.md}.

Re-run \texttt{fuzz}, the coverage and the goals whenever
\texttt{batch\_sender.py}, \texttt{recovery.py}, \texttt{display\_connection.py},
\texttt{dial\_generation.py}, \texttt{reply\_sink.py},
\texttt{liveness\_probe.py}, \texttt{replicator.py}, \texttt{death\_charger.py} or
the display's ack in \texttt{render\_loop.py} change who may charge a death, from
what evidence, when, or what an ack means.

\end{document}
